\clearpage
\section{판정 절차와 장 요약}
\label{sec:radicals-solvable-workflow}

\begin{learninggoal}
주어진 다항식이 근호로 풀리는지 판정하는 재현 가능한 절차를 정리한다. 갈루아군의 정규열에서 실제 근호탑을 읽는 방법과 삼차·사차 공식의 구조를 하나의 흐름으로 연결한다.
\end{learninggoal}

\subsection{근호해법 존재 판정}

분리다항식 $f\in K[X]$에 대해 다음 순서를 사용할 수 있다.

\begin{enumerate}
  \item \textbf{분리성 확인}: $\gcd(f,f')=1$인지 확인한다. 양의 표수에서는 비분리부분과 아르틴--슈라이어부분을 구분한다.
  \item \textbf{기약인수와 궤도 확인}: 기약성이 있으면 갈루아군은 근 집합에 추이적으로 작용한다. 인수분해는 궤도분해를 준다.
  \item \textbf{갈루아군 제한}: 판별식, 분해식, 좋은 소수의 인수분해형, 실근의 개수와 복소켤레를 사용해 $G\le S_n$을 좁힌다.
  \item \textbf{가해성 판정}: 알려진 군표, 도함열, 정규부분군과 몫군을 사용하여 $G$가 가해인지 판정한다.
  \item \textbf{정규열 선택}: 가능하면 인자군이 소수차 순환군인
  \[
    G=G_0\trianglerighteq\cdots\trianglerighteq G_r=1
  \]
  을 택한다.
  \item \textbf{고정체로 번역}: 갈루아 대응으로
  \[
    K=E_0\subseteq\cdots\subseteq E_r=L
  \]
  을 얻는다.
  \item \textbf{필요한 단위근 준비}: $E_{i+1}/E_i$의 차수가 $\ell\ne\charac K$이면 원시 $\ell$차 단위근을 먼저 넣는다.
  \item \textbf{근호 하나씩 첨가}: Kummer 또는 아르틴--슈라이어 정리로 각 층을
  \[
    \alpha_i^\ell\in E_i
    \quad\text{또는}\quad
    \alpha_i^p-\alpha_i\in E_i
  \]
  로 표현한다.
  \item \textbf{복원과 검산}: 얻은 원소들로 원래 근을 복원하고, 최소다항식·노름·대각합·Vieta 관계로 검산한다.
\end{enumerate}

\begin{pitfall}
갈루아군이 가해라는 사실은 근호공식의 \emph{존재}를 보장하지만, 계산하기 쉬운 공식을 자동으로 제공하지는 않는다. 정리의 증명에서 Kummer 생성원은 Hilbert 정리 90을 통해 존재하지만, 구체적인 생성원을 찾는 일은 별도의 분해식 계산을 요구한다.
\end{pitfall}

\subsection{자주 나타나는 군의 판정표}

\begin{center}
\renewcommand{\arraystretch}{1.22}
\begin{tabular}{p{0.18\textwidth}p{0.20\textwidth}p{0.22\textwidth}p{0.28\textwidth}}
\toprule
군 & 가해 여부 & 정규열의 예 & 방정식에서의 의미 \\
\midrule
$C_n$ & 가해 & 순환군의 소수차 세분 & 단위근을 준비하면 반복 근호첨가 \\
$V_4$ & 가해 & $V_4\trianglerighteq C_2\trianglerighteq1$ & 두 독립 제곱근 \\
$S_3$ & 가해 & $S_3\trianglerighteq C_3\trianglerighteq1$ & 판별식 제곱근 뒤 세제곱근 \\
$D_4$ & 가해 & $D_4\trianglerighteq C_4\trianglerighteq C_2\trianglerighteq1$ & 사차식의 반복 제곱근 층 \\
$A_4$ & 가해 & $A_4\trianglerighteq V_4\trianglerighteq C_2\trianglerighteq1$ & 삼차층과 제곱근 층 \\
$S_4$ & 가해 & $S_4\trianglerighteq A_4\trianglerighteq V_4\trianglerighteq C_2\trianglerighteq1$ & 일반 사차공식 \\
$D_5$ & 가해 & $D_5\trianglerighteq C_5\trianglerighteq1$ & 근호로 풀리는 오차식 \\
$F_{20}$ & 가해 & $F_{20}\trianglerighteq C_5\trianglerighteq1$ & Kummer형 오차식 \\
$A_5$ & 비가해 & $[A_5,A_5]=A_5$ & 근호해법 불가능 \\
$S_5$ & 비가해 & $S_5' =A_5$, $A_5'=A_5$ & 일반 오차식의 장애 \\
\bottomrule
\end{tabular}
\end{center}

\subsection{네 가지 대표 흐름}

\paragraph{기약 삼차식, 갈루아군 $S_3$.}
\[
  S_3\trianglerighteq A_3\trianglerighteq1
\]
에 따라 먼저 $\sqrt\Delta$를 첨가한다. 원시 세제곱근을 준비하면 남은 순환 삼차확장은 Lagrange 분해식의 세제곱근으로 생성된다.

\paragraph{기약 사차식, 갈루아군 $S_4$.}
세 쌍분할에 대한 몫작용 $S_4/V_4\cong S_3$을 삼차분해식으로 푼다. 그 뒤 $V_4$ 층은 $\sqrt{-z_1},\sqrt{-z_2},\sqrt{-z_3}$으로 해소된다.

\paragraph{기약 오차식, 갈루아군 $D_5$ 또는 $F_{20}$.}
두 군 모두 정규 $C_5$를 갖고 몫이 각각 $C_2$, $C_4$이다. 필요한 단위근을 첨가하면 $C_5$ 층은 오차근호 하나로 표현된다. 공식은 매우 복잡할 수 있지만 존재한다.

\paragraph{기약 오차식, 갈루아군 $A_5$ 또는 $S_5$.}
$A_5$가 완전군
\[
  [A_5,A_5]=A_5
\]
이므로 도함열이 멈추지 않는다. 따라서 어떠한 유한 근호탑도 분해체를 포함할 수 없다.

\subsection{무엇이 증명되었고 무엇이 남는가}

이 장에서 증명한 것은 다음과 같다.
\begin{itemize}
  \item 근호해법의 존재와 갈루아군의 가해성은 동치이다.
  \item 차수 $4$ 이하의 모든 분리방정식은 근호로 풀린다.
  \item 일반 $n$차방정식은 $n\ge5$에서 근호로 풀리지 않는다.
  \item 소수차 기약방정식에서는 실근과 비실근의 배치만으로 비가해성을 판정할 수 있는 강한 기준이 있다.
  \item Cardano와 사차공식의 근호 차수는 $S_3,S_4$의 정규열에서 나온다.
\end{itemize}

다음 장에서는 이 정리를 일반 오차방정식에 적용하기 위해 $A_5$의 단순성과 $S_n$의 구조를 더 깊게 분석한다. 목표는 ``일반 오차방정식은 왜 근호로 풀리지 않는가''를 완전한 군론적 증명으로 마무리하는 것이다.

\begin{chapterreview}
\begin{itemize}
  \item 교환자부분군 $[G,G]$은 $G$의 비가환성을 모은 가장 작은 정규부분군이다.
  \item $G$가 가해일 필요충분조건은 도함열이 유한번 뒤 $1$이 되는 것이다.
  \item 유한 가해군은 소수차 순환 인자군을 갖는 정규열을 가진다.
  \item 가해성은 부분군, 몫군, 가해 핵과 가해 몫의 확대에서 보존된다.
  \item $S_n$은 $n\le4$에서 가해이고 $n\ge5$에서 비가해이다.
  \item 근호탑에는 단위근, 표수와 서로소인 지수의 순수근호, 양의 표수의 아르틴--슈라이어 근이 허용된다.
  \item 유한확장이 근호로 풀릴 필요충분조건은 가해확장인 것이다.
  \item 분리다항식은 갈루아군이 가해일 때 정확히 근호로 풀린다.
  \item 가해인 소수차 추이군은 $\operatorname{AGL}(1,p)$ 안에 들어간다.
  \item 기약 소수차다항식이 두 실근 이상과 비실근을 가지면 근호로 풀리지 않는다.
  \item Cardano 공식의 제곱근과 세제곱근은 $S_3\trianglerighteq A_3\trianglerighteq1$에서 나온다.
  \item 사차공식의 삼차분해식과 세 제곱근은 $S_4/V_4$와 $V_4$ 층을 반영한다.
\end{itemize}
\end{chapterreview}

\subsection*{복습 카드}

\begin{itemize}
  \item \textbf{교환자부분군은?} 모든 $aba^{-1}b^{-1}$가 생성하는 부분군.
  \item \textbf{가해군의 도함열 판정은?} 어떤 $r$에 대해 $G^{(r)}=1$.
  \item \textbf{유한 가해군의 원자 인자군은?} 소수차 순환군.
  \item \textbf{근호해법의 갈루아 판정은?} 분해체의 갈루아군이 가해군.
  \item \textbf{$S_n$은 언제 가해인가?} 정확히 $n\le4$.
  \item \textbf{소수차 가해 추이군의 형태는?} $\operatorname{AGL}(1,p)$의 부분군.
  \item \textbf{삼차식의 첫 근호는?} 판별식의 제곱근.
  \item \textbf{사차식이 먼저 푸는 보조방정식은?} 세 쌍분할을 기록하는 삼차분해식.
\end{itemize}
